\begin{table}[!htbp]
\centering
\caption{The earlier-start demand-ratio window}\label{tab:long_demand}
\small
\setlength{\tabcolsep}{5pt}
\begin{tabular}{llrrrr}
\toprule
Type & Period & Short $N$ & Short $M$ & Long $N$ & Long $M$ \\
\midrule
Addition & 2007--11 & 8 & 57.87 & 8 & 60.83 \\
Addition & 2012--16 & 27 & -3.76 & 27 & -1.82 \\
Addition & 2017--21 & 58 & 2.97 & 58 & 4.24 \\
Addition & 2022--26 & 78 & 2.35 & 75 & 6.66 \\
Deletion & 2007--11 & 7 & -0.72 & 7 & 12.98 \\
Deletion & 2012--16 & 25 & 4.36 & 25 & 8.49 \\
Deletion & 2017--21 & 40 & 0.20 & 40 & -0.57 \\
Deletion & 2022--26 & 68 & 1.41 & 67 & 2.44 \\
\bottomrule
\end{tabular}
\begin{minipage}{\linewidth}\footnotesize\vspace{5pt}
Short is AD-1 to ED+1; Long is AD-20 to ED+1, the earlier-start counterpart in the anchor study. Each ratio uses identical events in its own numerator and denominator, but the short and long samples can differ. These are descriptive point estimates; the table does not claim well-identified early multipliers. The ED+20 analysis is a separate extension.
\end{minipage}
\end{table}
